\section{Лекция 1: Введение: от битов к пекетам}

Отчётеость:
\begin{itemize}
   \item курсовая (+ лабы~-- зачёт)
   \item экзамен
\end{itemize}

\textit{10 лекции и 4 лабы на весь курс.}

\subsubsection{Литература}

Компьютерные сети. Принципы, технологии протоколы

\subsection{Проблематика телекоммуникации}

\begin{itemize}
   \item Физическая передача данных по линиям связи (как нам передать сигнал, как осуществить физическую передачу данных).
   \item Объединения большого числа компонентов: (1) адресация, (2) безопасность, (3) организация общения большкго числа компонентов (маршрутизация).
   \item Обеспечение качества передачи: (1) скорость передачи, (2) расстояние, на которое можно передать (низкая задержка -> чем меньше заержа, тем лучше).
   \item Обеспечение безопасности: (1) доступность, (2) целостность, (3) конфиденциальность.
   \item Экономичекая эффективность \textit{(побеждает зачастую всё остальное, где сильно вложились в инфру, до сих пор исп. старые линии)}.
   \item Администрирование.
   \item Совместимость (все, кто делает оборудку, должны понимать, что к ним может прийти любой абонент).
\end{itemize}

\textit{Всё это ваимосвязано.}

\subsection{Интернет - система очередей. Коммутация каналов}

Интернет - живая, аморфная среда, каждая со своей очередью. Загрузки носят пиковый характер (трафик непостоянный). Схема соединения каждый с каждым (но! очень много соединений, и большинство из нх не импользуется).

Второй сценарий – коммутация каналов. Классический пример – телефон.
В Сетях коммутации линии бывают заняты, но такие сети идеально подходят для непрерывного трафика \textit{(всё время передаётся всё, что происходит)}.

Большитво трафика (ранее) тексторые нерешулярные каналы, поэтому используются сети коммутации пакетов.

\begin{itemize}
    \item Информация делится на части с заданной стактурой – пакеты.
    \item Чтобы решить проблему конфликта скорости, используется буфер (очередь пакетов).
    \item Каждый пакет обрабатывается независмио.
    \item Если скорость поступления больше скорости передачи, очередь начинает накапливаться.
    \item Чем больше очередь, тем больше задержка (задержка бывает очень низкой, а вечером, напимер, очень высокой).
    \item Переполнение очереди – потеря пакетов.
\end{itemize}

\[
\text{Задержка} = \text{Обработка} + \text{Время ожидания в очереди} + \text{Передача}
\]

\[
\text{Размер очереди} = \text{Пропускная способность} * \text{Среднее время ожидания пакета в очереди}
\]

Способы освобождений очереди:
\begin{itemize}
    \item Drop-tail (FIFO) \textit{(не всегла честно)}
    \item Случайные пакеты (RED, Blue) \textit{(кто уже стоит не в конце тоже может быть отброшен)}
    \item Отметки ECN \textit{(если близко к переполнению, каждый пакет помечается флагом – скорее информативные метки)}
\end{itemize}

\subsection{Иерархаческая органзация сетя. Декомпозиция}

Общий подход:
\begin{itemize}
    \item Разбиение на простые подзадачи.
    \item Чёткое определение функции.
    \item Выходные и входные интерфейсы.
    \item Координация.
\end{itemize}

Разбиваем на маленькие подзадачи. Работаем как конвейер. Общение происходит на уровням (как идёт информация каждому уровню без разницы) – на этом принципе построено всё общение компьютерных сетей.

Правила взаимодействия на уровнях:
\begin{itemize}
   \item Правила взаимодействия на уровнях – протоколы.
   \item Правила взаимодействия между уровнями – интерфейсы.
\end{itemize}

\textit{На экзамене знать семиуровнению модель OSI и принципы взаимодействия между уровнями.}

\begin{itemize}
    \item Сигнал в проводе – физический уровень.
    \item Коммутатор – физический и канальный уровни.
    \item 4-7 уровни – прокси сервер (замещает для отсальной сети вас).
\end{itemize}

\subsection{Физический уровень}

\begin{itemize}
    \item физическая среда
    \item дальность передачи
    \item тип кодирования
    \item синхронизация
    \item уровни сигналов
    \item разъёмы (независимо от того, что произвёл порт, всё будет работать)
\end{itemize}

Все функции физического уровня – аппарвтные.

Для физического уровня существует вероятность ошибки.

10 BaseT - UTP кат 3, 100 Ом, 100 м, 10 Мбит/с

\subsection{Канальный уровень}

Два подуровня:
\begin{itemize}
    \item MAC
    \item LLC (канальный уровень)
\end{itemize}

Ethernet, Token Ring, FDDI (типовый протоколы)

Кто говорит – канальный уровень.

\subsection{Сетевой уровень}

Основная обязанность – маршрутизация. Сетевой уровень отвечает за маршрутизацию в контексте вопроса где? (канальный – что? MAC адрес уникальный)

IP-адрес позваляет переать информацию до конкретного аобнента.

\begin{itemize}
    \item Внутри сети – канальный.
    \item Между – сетевой.
\end{itemize}

Реализуется программно. Функцианирует на уровне ядра ОС.

Протокол: IP

\subsection{Транспортный уровень}

В промежуточных узлах не обрабатывается. Отвечает за качество передачи данных. Транспортный уровень позволяет, например, открывать два браузера на одном компьютере (мультиплексирование).

Протокол: TCP/UDP, QUIC

\subsection{Сеансовый уровень}

Отвечает за управление диалогом. Задача установления и разрыва соединений. Синхронизация. Отдельно сеансового уровня практически не существует.

\subsection{Представительский уровень}

Задача унификации формы представления информации (собирать последовательностти правильно, чтобы переданное равнялось принятому).

Задача шифрования (протокол SSL)

\subsection{Прикладной уровень}

Уровень, на котором работают приложения. Большое число протоколов (HTTP, FTP, SMTP, SMTP). Клиент-серверная модель.

Любая модель декодирования данных – протокол.

\subsection{Инкапсуляия}

Каждый уровень добавляет свой заголовок к данным, пока не дойдём до физического уровня. В идеальной модели уровни не заглядаывют в заголовки других уровней, в жизни не так.

Чем меньше данных передаём, тем больше накладные расходы на передачу. Поэтому увеличение производительности сети – увеличение объёма данных.

На каждом уровне в заголовке существует информация, которая позволяет определить, откуда поступила информация. Позволяет делать мультиплексирование. \textit{(в реальности бывают нарушения последовательности)}

\subsection{Контроль передачи}

На канальном уровне – цикличекая контролная сумма (CRC-32). Сетевой – контроль чётности.

\subsection{Полоса и линия связи}

Полоса пропускания:

\[
\text{отношение апмлитуд} = \frac{A_{\text{выход}}}{A_{\text{вход}}} \geq \frac{1}{2}
\]

\subsection{Характеристика канала и линии связи}

Задержка:

\[
\text{задержка} ~ \frac{1}{2} * \text{RTT}
\]

Пропускная способность – фактическая достигнутая скорость передачи за единицу времени. Для 10 Мбит/c макс размер пакета: 9.76 Мбит/c

Bandwith-Delay Product (BDP):

\[
\text{BDP} = \text{пропускная способность} * \text{задержка}
\]

Макс объём данных, который может одновременно находиться <<в пути>>.

\section{Лекция 2}
\subsection{Канал (звено передачи данных)}

\begin{itemize}
    \item Пропсукная способность.
    \item Между – сетевой.
\end{itemize}

\begin{itemize}
    \item Доставка кадра.
    \item Адресация.
    \item Прообразование из битов в кадры.
    \item Синхронизация кадров.
    \item Обработка ошибок.
    \item Управление потоком.
\end{itemize}

\subsection{Области действия и функции DLL}

\begin{itemize}
    \item Локальные сетих: дооставка кадра между двумя узлами.
    \item Глобальные сети: доставка кадры между соседними узлами, соединёнными индивидуальной ЛС.
\end{itemize}

\begin{itemize}
    \item Синхронизация байтов, битов, кадров.
\end{itemize}

\subsection{Инкаспюляуия пакетов в кадры}

Заголовок: адрес назначения, адрес источника, тип, данные. Адрес назначения идёт первым, чтобы все слышали адрес назначения (прочитав первые 6 байт и понимаем, надо ли слушать информацию)

\subsection{Бит-ориентированные протоколы}

\textbf{Флаги}

Сначала идёт синхропроследовательнсть, после идёт флаг 01111110, обозначаем конец кадра 01111110. А внутри вставляем нолики, чтобы не было подряд 6 единиц. После 5 единиц вставляем 0 (бит-стаффинг).

Особенность: работает с битами, на уровне выше уже чёткое количесово байтов, поле длины у протокола отсутствует – есть начало и конец.

\textbf{Длина}

Даём преамбулу, потом начало кадро, потом заголовок, потом длина, потом данные. Так работает Ethernet.

\textbf{Использование запрещённых символов}

Преамбула, стартовый ограничитель, содержание кадра, стоповый ограничитель.
Используем стгланы, которых не может быть в нормальной передаче.

\subsection{Обработка ошибок}

Принцип: передача избыточной информации. Передаём степень мнгочлена $G(x) < \text{длины кадра}$, далее смотрим чётность.

\subsection{Особенности ЛС}

\begin{itemize}
    \item Типовая топология: шина, кольцо, звезда.
    \item Один маршрут (не должно быть петель).
    \item Разделяемая среда.
\end{itemize}

Достоинства: Простота и дешевизна.

\subsection{Подуровень LLC}

1 - без установления соединения и подтверж

\subsection{MACS}

Доступ к разлделяемой среде: случайный (Ethernet), детерминированный (Token Ring).

\subsection{MAC-адрес}

Физический адрес 6 байт, где саршие 3 байта – ID произвоителя (MAC в теории может быть одинаковым).

Адреса: инд (старший бит старшего байта = 0), групп (старший бит старшего байта = 1), Широковещательный (0xffffffffffff).

% TODO: дальше вставить скриншоты из лекции
